\chapter{Bad Breath}
\index{Bad Breath}

\section*{Definition and Overview}
Bad breath, also called halitosis, is a persistent unpleasant odour from the mouth. Nearly everyone has temporary bad breath at some time, such as ``morning breath'' or after eating strong foods, but halitosis usually refers to an ongoing problem that causes embarrassment and affects social life.

In the great majority of cases the source is the mouth itself: bacteria that live on the tongue and between the teeth break down food particles and dead cells, releasing smelly sulphur compounds. Sometimes the cause is elsewhere, such as the nose, throat, or lungs, and in a small number of cases it points to an underlying medical condition. Because most causes are treatable, persistent bad breath is worth investigating rather than living with.

\section*{Epidemiology}
Bad breath is common, with studies suggesting that around 1 in 4 people has persistent halitosis at some point. It becomes more common with age, and is more frequent in people with gum disease, poor oral hygiene, a dry mouth, and in smokers and people who drink alcohol heavily. Many people who worry about bad breath do not actually have it, and a smaller number have persistent concern about a non-existent odour (a condition called halitophobia). Bad breath is also a leading cause of dental visits after tooth decay.

\section*{Aetiology and Pathophysiology}
The odour is produced by bacteria breaking down proteins into volatile sulphur compounds. Causes include:

\begin{itemize}
    \item \textbf{Poor oral hygiene:} food trapped between teeth, plaque, and a heavily coated tongue
    \item \textbf{Gum disease (gingivitis and periodontitis):} inflamed gums with pockets where bacteria collect
    \item \textbf{Dental decay, broken fillings, and abscesses}
    \item \textbf{Tonsil stones (tonsilloliths):} small hardened deposits in the tonsils with a strong odour
    \item \textbf{Dry mouth (xerostomia):} reduced saliva from medications, mouth breathing, or conditions such as Sj\"ogren's syndrome
    \item \textbf{Nose and throat:} post-nasal drip, sinusitis, and throat infections
    \item \textbf{Foods and habits:} garlic, onions, spices, coffee, alcohol, and tobacco
    \item \textbf{Less common systemic causes:} poorly controlled diabetes (fruity or acetone smell), liver or kidney failure, reflux, and some infections
\end{itemize}

\section*{Clinical Features}
\begin{itemize}
    \item An unpleasant or persistent odour from the mouth, often noticed by others rather than the person themselves
    \item A bad or metallic taste in the mouth
    \item A white or coated tongue
    \item Dry mouth, thick saliva, or a ``sticky'' feeling
    \item Associated dental problems: bleeding gums, gum tenderness, tooth pain, or loose teeth
    \item Frequent throat clearing or a sensation of something in the throat (tonsil stones)
    \item A persistent acid or bitter taste, or symptoms of reflux
\end{itemize}

\section*{Differential Diagnosis}
\begin{itemize}
    \item Oral causes (poor hygiene, gum disease, decay, tongue coating) versus nasal and sinus causes versus tonsil stones
    \item Dry mouth from medications versus a medical cause of reduced saliva
    \item Reflux (gastro-oesophageal reflux disease) as an uncommon contributor
    \item Systemic disease: diabetes, liver disease, or kidney disease with characteristic odours
    \item Halitophobia -- concern about bad breath that is not present, or that others cannot detect
\end{itemize}

\section*{When to Seek Medical Care}
Bad breath that persists despite good brushing, flossing, and tongue cleaning should be assessed. Start with a dentist, because the cause is usually dental.

See a dentist if:
\begin{itemize}
    \item Bad breath continues despite good oral hygiene
    \item You have bleeding, tender, or receding gums, tooth pain, or loose teeth
    \item You have not had a dental check-up for some time
\end{itemize}

See a doctor if bad breath is accompanied by:
\begin{itemize}
    \item Persistent acid reflux or heartburn
    \item Nasal congestion, post-nasal drip, or throat symptoms
    \item Excessive thirst and weight loss (possible diabetes)
    \item An unusual smell (fruity, fishy, or ammoniac) with feeling unwell
    \item Dry mouth, a burning mouth, or a mouth sore that will not heal
\end{itemize}

\section*{Diagnosis and Investigations}
\begin{itemize}
    \item \textbf{Dental examination:} a full check for decay, gum disease, and defective fillings
    \item \textbf{Tongue assessment:} looking for a heavy coating that is a common source of odour
    \item \textbf{Odour assessment:} sniffing tests or a halimeter, which measures sulphur compounds in the breath
    \item \textbf{ENT (ear, nose and throat) review:} for tonsil stones, sinusitis, or post-nasal drip
    \item \textbf{Further tests:} blood sugar tests if diabetes is suspected, or tests for reflux, liver, or kidney problems in appropriate cases
    \item \textbf{Dental X-rays:} to find decay or abscesses that are not visible
\end{itemize}

\section*{Management}
\begin{itemize}
    \item \textbf{Professional dental care:} scale and polish, treatment of gum disease and decay, and repair of broken fillings
    \item \textbf{Tongue cleaning:} gently brushing or scraping the back of the tongue each day
    \item \textbf{Oral hygiene:} brush twice daily with fluoride toothpaste and clean between the teeth with floss or interdental brushes
    \item \textbf{Mouthwashes:} an antibacterial or alcohol-free rinse, used short-term as advised
    \item \textbf{Dry mouth measures:} sip water, use sugar-free gum, and review medicines that reduce saliva
    \item \textbf{Treating specific causes:} removal of tonsil stones, treatment of sinusitis, management of reflux, or diabetes care
    \item \textbf{Lifestyle changes:} stop smoking, limit alcohol and strong foods, and stay hydrated
\end{itemize}

\section*{Complications}
\begin{itemize}
    \item Social and psychological effects: embarrassment, anxiety, and avoidance of close contact
    \item Halitophobia, where worry about breath becomes a persistent anxiety or obsession
    \item Untreated gum disease progressing to tooth loss
    \item An underlying condition (such as diabetes or sinus disease) remaining undiagnosed if the smell is not investigated
    \item Masking the odour with mints or mouthwash can delay treatment of the real cause
\end{itemize}

\section*{Prognosis and Outcomes}
Most cases of bad breath improve greatly, or resolve completely, with good oral hygiene and treatment of the underlying cause. Where gum disease, decay, or tonsil stones are responsible, professional treatment usually makes a lasting difference. A dry mouth or ongoing medical conditions such as reflux or diabetes may need ongoing management, and the smell can return if these are not controlled. In most people, consistent self-care plus regular dental visits keeps the problem manageable.

\section*{Prevention and Self-Care}
\begin{itemize}
    \item Brush teeth twice a day and clean between the teeth daily
    \item Brush or scrape the tongue each day
    \item Visit the dentist regularly, usually once or twice a year
    \item Drink plenty of water and avoid a dry mouth
    \item Use sugar-free gum after meals if you cannot brush
    \item Avoid tobacco, and limit alcohol, coffee, and strong-smelling foods
    \item Treat reflux, nasal congestion, and other contributing conditions
    \item If the problem persists despite good hygiene, seek professional advice rather than relying on mints and mouthwash
\end{itemize}

\section*{Sources and Further Reading}
The following sources provide reliable, up-to-date information on bad breath and oral health.
\begin{itemize}
    \item NHS. \textit{Bad breath (halitosis).} https://www.nhs.uk/conditions/
    \item Mayo Clinic. \textit{Bad breath.} https://www.mayoclinic.org
    \item MSD Manual. \textit{Halitosis.} https://www.msdmanuals.com
    \item MedlinePlus. \textit{Bad Breath.} https://medlineplus.gov
    \item Centers for Disease Control and Prevention. \textit{Oral Health.} https://www.cdc.gov
\end{itemize}
